\documentclass[english,twoside, openright, hidelinks, 11 pt, class=report,crop=false]{standalone}
\input{../../preamb}
\input{../bm}

\usepackage{xr}
\externaldocument{../MB_bm}
\begin{document}

\footnotesize

\subsection*{\faskap \ref{Talavare}}
\opr{opgtelverdi}
\abch{
\item 22 
\item 13
\item 36
}

\opr{opgtelverditl}
\abch{
\item 17
\item 11
\item 29
}

\opr{opgtelverdidestl}
\abch{
\item 1,7
\item 2,3
}

\opr{opgtelverdides}
\abch{
\item 13,9
\item 32,8
\item 0,7
\item 2,4
}

\opr{opgeininndelta} b)

\opr{opgeininndeltc} c)

\opr{opg100inndelta} a)

\opr{opg100inndelta} c)

\opr{opgtellargestsmallest}
\abch{\item 871 \item 178}

\opr{opgkord}
$A=(1, 3)$, $B=(2, 4)$, $C=(6, 1)$, $D=(5, 5)$

\opr{opgkord}
\includegraphics[scale=0.25]{\figp{opgkord_lf}}

\subsection*{\faskap \ref{Rekneartane}}
\opr{opgsumtotall}
\mulansw
\abch{
\item $ 4=1+3 $
\item $ 5=2+3 $
\item $ 6=2+4 $
\item $ 7=1+6 $
\item $ 8=3+5 $
\item $ 9=2+7 $
}

\opr{opgsumtretall}
\mulansw
\abch{
	\item $ 5 =2+2+1 $ 
	\item $ 6 = 1+3+2 $
	\item $ 7 = 3+2+2 $
	\item $ 8 = 1+1+6 $
	\item $ 9 = 3+3+3 $
}

\opr{opgtiervenn} \vspace{-5pt}

\abcn{
\item
	\abch{
		\item 8
		\item 7 
		\item 6
		\item 5
	}
\item  Fordi de er svarene i oppgave 1a), 1b) og 1c), og addisjon er kommutativ. 
}

\opr{opgsumparodd}
\abch{
\item 1)
\item 1)
\item 2)
}

\opr{opgdifto} \mulansw
\abch{
\item $ 2=7-5 $
\item $ 3=10-7 $
\item $ 4=5-1 $
\item $ 5=10-5 $
\item $ 6 = 9-3 $
\item $ 7 = 9-2 $
\item $ 8=10-2 $
}

\opr{opgdifparodd}
\abch{
	\item 1)
	\item 1)
	\item 2)
}

\opr{opgaddifomv}
\abch{
	\item 8 og 13
	\item 42 og 19
	\item 762 og 141
	\item 15 og 9
	\item 83 og 25
	\item 904 og 352
	\item 3796 og 52
}

\opr{opgadsomgang}
\abch{
	\item $ 2+2+2+2=2\cdot 4 = 4+4 $
	\item $ 3+3+3+3+3+3= 3\cdot 6=6+6+6 $
	\item $ 4+4=4\cdot2 = 2+2+2+2 $ 
	\item $ 5+5+5+5+5+5+5+5+5+5 = 5\cdot10=10+10+10+10+10 $
	\item $ 6+6+6+6+6=6\cdot5 = 5+5+5+5+5+5 $
	\item $ 7+7+7+7=7\cdot4=4+4+4+4+4+4+4 $
}

\opr{opggangrut}
\abch{
	\item \includegraphics{\figp{opggangruta}}
	\item \includegraphics{\figp{opggangrutb}}
	\item \includegraphics{\figp{opggangrutc}}
	\item \includegraphics{\figp{opggangrutd}}
	\item \includegraphics{\figp{opggangrute}}
}

\opr{opggangdivomv}
\abch{
\item 5 og 9
\item 8 og 23
\item 56 og 42
\item 21 og 3
\item 216 og 9
}

\opr{opggangparod} 
\abch{
	\item partall
	\item partall, 0
	\item oddetall, 5
}

\subsection*{\faskap \ref{Faktogrek}}

\opr{opggangutv}
\abch{
	\item 34
	\item 177
	\item 100
	\item 664
	\item 2943
}

\opr{opgfaktto} \mulansw
\abch{
	\item $ 5\cdot20 $
	\item $ 3\cdot10 $
	\item $ 2\cdot 20 $
	\item $ 2\cdot 35 $
	\item $ 7\cdot 6 $
	\ \ \item $ 8\cdot 4 $
	\item $ 42\cdot2 $
	\item $ 3\cdot30 $
}

\opr{opgprimfakt}
\abch{
	\item $ 2\cdot 2\cdot2\cdot2\cdot 3 $ 
	\item $2\cdot 2\cdot3 \cdot5 $
	\item $ 2\cdot 5 $
} 

\opr{opgfakttre}
\mulansw
\abch{ 
	\item $ 5\cdot20 $, $ 10\cdot10 $, $ 5\cdot2\cdot20 $, 
	\item $ 3\cdot10 $, $ 6\cdot5 $, $ 3\cdot2\cdot5 $
	\item $ 2\cdot 20 $, $ 4\cdot 10 $, $ 2\cdot2\cdot10 $
	\item $ 2\cdot 35 $, $ 14\cdot 5 $, $ 2\cdot7\cdot5 $
	\item $ 7\cdot 6 $, $ 7\cdot2\cdot3 $, $ 21\cdot 2 $
	\item $ 8\cdot 4 $, $ 16\cdot 2 $, $ 8\cdot \cdot2\cdot2 $
	\item $ 42\cdot2 $, $21\cdot4$, $7\cdot6\cdot2$
	\item $ 3\cdot30 $, $ 9\cdot10 $, $ 3\cdot3\cdot10 $
}

\grubr{opgforklprimfakt}

\grubr{opggangtoledd}
\abch{
\item \selos
\item 1715
}

\grubr{opgoddgangodd}

\grubr{opgperf} 28

\grubr{opgfindprime} \selos



\subsection*{\faskap \ref{Brok}}

\opr{opgbrverdidelelig}
\abch{
\item 6
\item 5
\item 2
\item 7
\item 9
\item 4
}

\opr{opgbrverdiikkedelelig}
\abch{
\item 0,5
\item 0,25
\item 0,2
\item 0,75
\item 0,4
\item 0,6
\item 0,8
\item 1,5
\item 0,33...
\item 2,5
\item 0,833...
\item 1,4
\item 2,75
\item 0,7
}

\opr{opgbrfigs}
\abch{
	\item $\frac{2}{3}$
	\item $\frac{5}{9}$
	\item $\frac{3}{7}$
	\item $\frac{7}{10}$
}

\opr{opgbrtallin}
\abch{
\item $ \frac{3}{4} $
\item $ \frac{2}{7} $
\item $ \frac{2}{5} $
}

\opr{opgbrtallin2}
\abch{
\item $ \frac{14}{3} $
\item $ \frac{13}{2} $
\item $ \frac{11}{5} $
}

\opr{opgbrutvid}
\abch{
\item $ \frac{20}{6} $
\item $ \frac{9}{12} $
\item $ \frac{12}{28} $
\item $ \frac{45}{40} $
\item $ \frac{54}{30} $
\item $ \frac{77}{28} $
}

\opr{opgbrutvidtil}
\abch{
\item $ \frac{9}{4} $
\item $ \frac{9}{5} $
\item $ \frac{2}{7} $
}

\opr{opgbrsjekk}
\abch{
	\item $\frac{4}{3}>\frac{6}{5}$
	\item $\frac{20}{50}=\frac{2}{5}$
	\item $\frac{7}{3}>\frac{9}{5}$
}

\opr{opgbrad}
\abch{
\item $ \frac{10}{3} $
\item $ \frac{14}{4} $
\item $ \frac{11}{6} $
\item $ \frac{10}{7} $
\item 1
} \\
\mers{Én av brøkene kan forkortes.}

\opr{opgbrsub}
\abch{
	\item $ \frac{1}{3} $
	\item $ \frac{2}{4} $
	\item $ \frac{9}{6} $
	\item 1
	\item 0
} \\
\mers{To av brøkene kan forkortes.}

\opr{opgbrad2}
\abch{
	\item $ \frac{22}{3} $
	\item $ \frac{8}{5} $
	\item $ \frac{17}{7} $
}

\opr{opgbradandsub}
\abch{
\item $ \frac{6}{5} $	
\item $ \frac{5}{7} $
\item 4
}


\opr{opgbrbrandin}
\abch{
	\item $\frac{17}{7}$
	\item $\frac{29}{4}$
}


\opr{opgbrad3}
\abch{
\item $ \frac{9}{10} $
\item $ \frac{73}{63} $
\item $ \frac{101}{24} $
\item $ \frac{73}{20} $
\item $ \frac{5}{6} $
}

\opr{opgbrsub2}
\abch{
\item $ \frac{1}{10} $
\item $ \frac{29}{36} $
\item $ \frac{71}{72} $
\item $ \frac{11}{20} $
\item $ \frac{5}{6} $
}

\opr{opgbradandsubmix}
\abch{
\item $ \frac{5}{12} $
\item $ \frac{157}{30} $
\item $ \frac{229}{56} $
}

\opr{opgbrgongheil}
\abch{
\item $ \frac{20}{3} $
\item $ \frac{40}{7} $
\item $ \frac{54}{10} $
\item $ \frac{80}{7} $
\item $ \frac{21}{2} $
\item $ \frac{28}{3} $
\item $ \frac{35}{3} $
\item $ \frac{30}{7} $
\item $ \frac{5}{11} $
\item $ \frac{72}{17} $ 
}

\opr{opgbrdelheil}
\abch{
\item $ \frac{4}{15} $
\item $ \frac{5}{56} $
\item $ \frac{9}{60} $
\item $ \frac{8}{70} $
\item $ \frac{3}{14} $
\item $ \frac{9}{110} $
\item $ \frac{1}{60} $
\item $ \frac{9}{290} $
\item $ \frac{8}{459} $
\item $ \frac{4}{158} $
}

\opr{opgbrdelheldel}
\abch{
\item $\frac{2}{3}$
\item $\frac{2}{5}$
\item $\frac{4}{3}$
\item $\frac{6}{13}$
}

\opr{opgbrgongbr}
\abch{
\item $ \frac{20}{27} $
\item $ \frac{7}{32} $
\item $ \frac{18}{21} $
\item $ \frac{60}{5} $
\item $ \frac{21}{10} $
\item $ \frac{10}{21} $
\item $ \frac{16}{21} $
\item $ \frac{80}{9} $
\item $ \frac{36}{35} $
\item $ \frac{35}{12} $
}

\opr{opgbrgongbr2}
\abch{
\item $ \frac{15}{40} $
\item $ \frac{153}{32} $
\item $ \frac{46}{32} $
\item $ \frac{21}{648} $
\item $ \frac{203}{328} $
}

\opr{kansfakt}
\abch{
\item $ \frac{4}{11} $
\item $ \frac{35}{8} $
\item $ \frac{1}{9} $
\item 4
}

\opr{opgbrforkortmfaktor}
\abch{
\item $ \frac{7}{4} $
\item $ \frac{3}{7} $
\item $ \frac{2}{3} $
\item $ \frac{8}{7} $
\item $ \frac{1}{2} $
\item $ \frac{7}{2} $
} \vsk


\opr{opgbrgongbrfakt}
\abch{
\item 49
\item 54
\item 70
\item 16
\item 30
\item 12
\item 25
\item 14
\item 7
\item 63
}

\opr{opgbrdelbr}
\abch{
\item $ \frac{14}{15} $
\item $ \frac{24}{45} $
\item $ \frac{30}{21} $
\item $ \frac{7}{20} $
\item $ \frac{66}{15} $
}

\opr{opgbrdelbr2}
\abch{
	\item $ \frac{4}{9} $
	\item $ \frac{15}{4} $
	\item $ \frac{14}{5} $
} \os

\grubr{opgfracval} 
\abch{
\item 2
\item 4
\item 5
\item 2
\item 4
\item 5
\item 3, 4
\item 2, 5
\item 3, 5
\item 4, 5	
} \os

\grubr{opgfracval2}
\abch{
	\item 2
	\item 4
	\item 5
	\item 4, 3
	\item 5, 2
	\item 5, 3
	\item 5, 4
} \os

\grubr{opgfracprim}
\abch{
\item $ \frac{403}{6732} $
\item $ \frac{269}{3150} $
}

\subsection*{\faskap \ref{Negtal}}

\opr{opgnegfyllinn}
\abch{
	\item 9 har retning mot høgre og lengde 9.
	\item 4 har retning mot høgre og lengde 4
	\item -3 har retning mot venstre og lengde 3
	\item 12 har retning mot høgre og lengde 12
	\item -11 har retning mot venstre og lengde 11
	\item -25 har retning mot venstre og lengde 25
}	

\opr{opgnegplas}\\
\includegraphics{\figp{opgneg1fas}}


\opr{opgnegogpos}
\abch{
	\item 9, 4 og 12.
	\item $ -3 $, $ -11 $ og $ -25 $.
}


\opr{opgnegad}
\abch{
	\item $ 1 $
	\item $ 8 $
	\item $ 6 $
	\item $ 0 $
	\item $ 3 $
	\ \ \ \item $ 1 $
	\ \ \ \item $ 10 $
	\item $ 5 $
}

\opr{opgnegad2}
\abch{
	\item $ -16 $
	\item $ -8 $
	\item $ -23 $
	\item $ 0 $
	\item $ -23 $
	\item $ -17 $
	\item $ -3 $
	\item $ 0 $
}

\opr{opgnegsub}
\abch{
	\item $ 15 $
	\item $ 17 $
	\item $ 12 $
	\item $ 14 $
	\item $ -12 $
	\ \ \ \item $ -19 $
	\ \ \ \item $ -12 $
	\item $ -13 $
}

\opr{opgnegsub2}
\abch{
	\item $ 22 $
	\item $ 22 $
	\item $ -17 $
	\item $ 14 $
	\item $ 15 $
	\item $ 13 $
	\item $ -13 $
	\item $ -12 $
}

\opr{opgnegsubadbrok}
\abch{
	\item $-\frac{5}{3}$
	\item $\frac{1}{5}$
	\item $-\frac{8}{11}$
	\item $\frac{3}{7}$
}

\opr{opgnegmult}
\abch{
	\item $ -12 $
	\item $ -50 $
	\item $ -63 $
	\item $ -24 $
	\item $ -56 $
	\item $ -27 $
	\item $ -12 $
	\item $ -40 $
	\item $ 21 $
	\item $ 25 $
	\item $ 12 $
	\item $ 72 $
}

\opr{opgnegdiv}
\abch{
	\item $ -4 $
	\item $ -6 $
	\item $ -5 $
	\item $ -4 $
	\item $ -8 $
	\item $ -9 $
	\item $ -5 $
	\item $ -5$
	\item $ 8 $
	\item $ 9$
	\item $ 5 $
}

\opr{opgkordneg}
\includegraphics[scale=0.25]{\figp{opgkordneg}}

\subsection*{\faskap \ref{Utrekning}}
\opr{sumuovg}
\abch{
	\item $ 96 $
	\item $ 87 $	
	\item $ 899 $
	\item $ 298 $
}

\opr{summovg}
\abch{
	\item $ 103 $
	\item $ 143 $	
	\item $ 1000 $
	\item $ 1010 $
}


\opr{difuovg}
\abch{
	\item $ 61 $
	\item $ 234 $	
	\item $ 128 $
	\item $ 651 $
}

\opr{difmovg}
\abch{
	\item $ 59 $
	\item $ 325 $	
	\item $ 523 $
	\item $ 87 $
}


\opr{gangetsif}
\abch{
	\item $ 36 $
	\item $ 112 $	
	\item $ 380 $
	\item $ 258 $
	\item $ 763 $
	\item $ 784 $
	\item $ 1917 $
}

\opr{gangtosif}
\abch{
	\item $ 348 $
	\item $ 2573 $
	\item $ 6916 $
	\item $ 1162 $
}

\opr{gangtretosif}
\abch{
	\item $ 29\,736 $
	\item $ 58\,183 $
	\item $ 61\,243 $
	\item $ 48\,977 $
}

\opr{gangdes1}
\abch{
	\item $ 135 $ og $ 13,5 $
	\item $ 47\,424 $ og $ 474,24 $
	\item $ 68\,672 $ og $ 68,672 $
	\item $ 569\,366 $ og $ 56,9366 $
}

\opr{gangdes2}
\abch{
	\item $ 411,5 $ 
	\item $ 66,57 $ 
	\item $ 38,08 $
}



\opr{divtoensif}
\abch{
	\item $ 49 $
	\item $ 29 $
	\item $ 23 $
	\item $ 17 $
	\item $ 12 $
} 

\opr{divtrensif}
\abch{
	\item $ 189 $
	\item $ 56 $
	\item $ 99 $
	\item $ 19 $
	\item $ 26 $
	\item $ 97 $
}



\opr{stndform1}
\abch{
	\item $ 9,8\cdot 10^4 $
	\item $ 1,67\cdot 10^8 $
	\item $ 4,819\cdot 10^3 $
	\item $ 2,1\cdot 10^1 $
	\item $ 9,13227\cdot 10^3 $
	\item $ 8,937\cdot 10^2 $
	\item $ 1,80021\cdot 10^4 $
	\item $ 3,024\cdot 10^2 $
}

\newpage
\opr{stndform2}
\abch{
	\item $ 2,7\cdot 10^{-2} $
	\item $ 1,901\cdot 10^{-4} $
	\item $ 3,2\cdot 10^{-1} $
	\item $ 2,0032\cdot 10^{-7} $
} 

\grubr{stnddroft}
\abch{
	\item $ 9\cdot 10^8 \cdot 7\cdot10^{-5} $
	\item $ 6,3\cdot 10^4 $
}


\subsection*{\faskap \ref{Geometri}}
\opr{opgomkr1}
\abch{
\item 14
\item 20
\item 24
}

\opr{opgarfirk}
\abch{
\item 2
\item 16
\item 27
}

\opr{opgarfinn}
\abch{
\item 2 og 8
\item 3 og 4
\item 3 og 6
}

\opr{opgarstorst}
\abch{
\item 81
\item 1) Et åpenbart eksempel er kvadratet fra oppgave a), som har bredde og høyde 9, og areal 81. 2) Bredde 15 og høyde 3, areal 45. 3) Bredde 12 og høyde 6, areal 72. \mulansw
}

\opr{opgartrek}
\abch{
\item 3
\item 10
\item 6 
\item 6
\item 10
\item 6
\item 4
\item 3
\item 28
}

\opr{opggeovolprism}
\abch{
\item 90. \mers{Grunnflatearealet kan også være 72 eller 80, avhengig av hvilken sideflate man velger ut som grunnflate.}
}

\opr{opgtr1} \selos 

\opr{opgsp1} \selos

\grubr{opgomkrogarparodd} \selos

\grubr{opgtrkkvadv} \selos

\grubr{opggeolikar} \selos

\subsection*{\faskap \ref{Algebra}}
\opr{opgalggjentad}
\abch{
	\item $ 3a $ \item $ 4a $ \item $ 7a $
	\item $ -2b $
	\item $ -5b $
	\item $ -3k $
}

\opr{opgalgadogsub}
\abch{
	\item $ a+b $
	\item $ a+2b $
	\item $ 9b-3a $
} \vsk

\opr{opgalgadogsub2}
\abch{
	\item $ 2b-5a+c $
	\item $ 3b-9a $
	\item $ 11b-3a $
}

\opr{opgalggong}
\abch{
	\item $ 7a+14 $
	\item $ 9b+27 $
	\item $ 8b-24c $
	\item $ -6a-10b $
	\item $ (9a+2) $
	\item $ (3b+8)a $
	\item $ 3ac-ab $
	\item $ 2a+6b+8c$
	\item $ 27b-9c+63a $
	\item $ 2c-6b-14a $
}

\opr{opgalgfaktoriser}
\abch{
	\item $ 2(a+b) $
	\item $ b(4a+5) $
	\item $ c(9b-1) $
	\item $ 2a(2c-1) $
}

\opr{opgalgkvad}

\opr{GV21D1opg10}
\abch{
	\item $ \frac{1}{4} $
	\item $ \frac{1}{2} $ når $ x=4 $ og $ 2 $ når $ x=-2 $.
}

\opr{eksu22opg2}
Uttrykkene gitt av a) og c) stemmer.

\opr{opgpot}
\abch{
	\item $ 3^4 $
	\item $ 5^2$
	\item $ 7^6 $
	\item $ a^3 $
	\item $ b^2 $
	\item $ (-c)^4 $ \mers{$ (-c)^4=c^4 $}
}

\opr{opgpotverdi}
\abch{
	\item $ 64 $
	\item $ 32 $
	\item $ 64 $
	\item $ -8$
	\item $ -243 $
	\item $ 256 $
} 

\opr{opgpotskrivtilpot}
\abch{
	\item $ 2^{16} $
	\item $ 3^{11} $
	\item $ 9^6 $
	\item $ 6^5 $
	\item $ 5^{-4} $
	\item $ 10^{11} $
	\item $ a^{16} $
	\item $ k^7 $
	\item $ x^3 $
	\item $ x$
	\item $ 1 $
	\item $ a^5\cdot b^{-3} $
}

\opr{opgrotter}
\abch{
	\item $ 5 $
	\item $ 10 $
	\item $ 12 $
	\item $ 3 $
	\item $ 9 $
	\item $ 10 $
}

\grubr{1Tpotens}
\abch{
\item $\frac{2}{3}$
\item $ \frac{1}{8}$ 
\item $\frac{1}{9}a^{-\frac{1}{2}}$
\item $\frac{7}{5}$
}

\grubr{1Tforenkle}
\abch{
\item $x-y$
\item $\frac{2(x+1)}{x-1}$
}

\grubr{1Tforenkle2}
\abch{
\item $\frac{x+1}{x-3}$
\item $\frac{x+2}{x-2}$
\item $\frac{3(x+1)}{x-1}$
\item $\frac{1}{x-3}$
}

\grubr{1Tstandardform}
\abch{
	\item $5\cdot10^9$
	\item $2\cdot10^7$
	\item $3,1\cdot10^16$
	\item $1,9\cdot10^{-4}$
}

\grubr{opgalglikbrok}
\selos

\grubr{opgalgtverrsum}
\selos
 
\subsection*{\faskap \ref{Likningar}}

\opr{opglig3ledd}
\abch{
\item $ x= 10$
\item $ x= 5$
\item $ x= 9$
\item $ x= 2$
\item $ x=3 $
\item $ x=4 $
\item $ x= 10$
\item $ x= 8$
\item $ x= 10$
}

\opr{opglig4ledd}
\abch{
\item $ x= 37$
\item $ x= 29$
\item $ x= 23$
\item $ x= 15$
\item $ x= 8$
\item $ x= 11$
\item $ x= 23$
\item $ x= 19$
\item $ x= 10$
\item $ x= 28$
}

\opr{opgligax1}

\abch{
\item $ x= 4$
\item $ x= 5$
\item $ x= 9$
\item $ x= 15$
}

\opr{opgligxovera1}
\abch{
\item $ x= 8$
\item $ x= 72$
\item $ x= 49$
\item $ x= 150$
}

\opr{opgligax2}
\abch{
\item $ x= 7$
\item $ x= 10 $
\item $ x= 6$
\item $ x= 9$
\item $ x= 9$
\item $ x= 8$
\item $ x= 5$
\item $ x=2 $
}

\opr{opgligvink} \selos

\opr{eksu22opg3} $ x=18 $

\opr{opgeqeqsys}
\abch{
\item $x=6$, $y=19$
\item $x=3$, $y=7$
\item $x=3$, $y=0$
\item $x=\frac{63}{5}$, $y=18$
}

\opr{1TV21D1opg1} $ x=1, y=-2 $

\grubr{1PV22D1opg6}
36

\grubr{opgeqwthfrac1}
\abch{
	\item $x=3$
	\item $x=2$
	\item $x=9$
	\item $x=11$
}

\grubr{opgeqwthfrac2}
\abch{
	\item $x=-\frac{21}{2}$
	\item $x=\frac{90}{23}$
	\item $x=27$
	\item $x=\frac{30}{11}$
	\item $x=-6$
	\item $x=\frac{11}{17}$
}



\grubr{opgrepdes} \selos

\subsection*{\faskap \ref{Funksjoner}}
\opr{opgfunfinnlin}
\abch{
	\item $ f(x)=2x+4 $
	\item 204
	\item $ x=10 $
}

\opr{opgfunkfem}
\abch{
	\item $ a(x)=5x^2 $
	\item $ x = 2000 $
	\item $ x =9 $
}

\opr{opgfunkbluegreen}
\abch{
	\item $ b(x)=x^2+2x $
	\item 440
	\item $ x= 8 $
}

\opr{EG22D1opg1}
22 trekanter og 10 firkanter.

\opr{opgfunkparodd}
La $ x $ være et positivt heltall.
\abch{
	\item $ p(n)=2n $
	\item $ o(n)=2n-1 $
}

\opr{opgfunkstigogkonst}

\abch{
	\item Stigningstall $ 5 $ og konstantledd $ 10 $.
	\item Stigningstall $ 3 $ og konstantledd $ -12 $.
	\item  Stigningstall $ -\frac{1}{7} $ og konstantledd $ -9 $.
	\item Stigningstall $ \frac{3}{2} $ og konstantledd $ -\frac{1}{4} $.
}

\opr{funkopgtegn} \selos


\opr{funkopgligset}
\abch{
	\item \eqref{eks3a} og \eqref{eks3b} gir hver for seg en ligning som beskriver en rett linje. Disse linjene kan også representeres ved $ f $ og $ g $ slik som de er definert.
	\item $ x=7 $ og $ y=2 $.
}

\opr{funkopgfinngraf1} $ f(x)=-8x+16 $

\opr{funkopgfinngraf2} $ f(x)=4x+3 $

\opr{funkopgfinngraf3} $ f(x)=x-3 $ og $ g(x)=\frac{1}{4}x+1 $

\grubr{opgfunkfindab}
\abch{\item $y=4x+10$ \item $7x-2$ \item $-7x+10$ \item $-\frac{1}{3}x-5$}

\grubr{1ph21opg5} 
\abch{\item $f(x)=-2x+3$, $g(x)=\frac{1}{2}-2$ \item $\frac{5}{4}$}

\grubr{opgfunkpar}
\abch{
\item \includegraphics[scale=0.25]{\figp{opgfunkpar_lf}}
\item $f$ og $g$ er parallelle. \selos
}

\grubr{opgfunkfindslope}
$g(x)=2x$

\grubr{opgfunkrect1}
\abch{
\item 3
\item $\frac{3}{2}$
\item 3
}

\grubr{opgfunkrect2}
$\frac{5}{2}$

\grubr{opgfunkrectangleandline}
$4$


\grubr{opgaddifparodd}
\selos

\subsection*{\faskap \ref{Geometri2}}

\opr{opggeofractrap}
$\frac{62}{135}$

\opr{opggeoareaomv} 
\abch{
	\item 3
	\item $9b$
	\item 10
}

\opr{opggeocirc}
\abch{
	\item $8\pi$
	\item $16\pi$
}

\opr{opggeocircomv}
\abch{
	\item 10
	\item 9
	\item $2k$
	\item $10a$
}
	
\opr{opggeocirkar}
\abch{
	\item $6\pi$ og $24\pi$
	\item $9\pi$ og $144\pi$
	\item 4
	\item 16
}

\opr{opggeocirkargeneral}
\abch{
	\item $a$
	\item $a^2$
}	

\opr{opggeovolkjegl}
\abch{
	\item $ 100\pi $
	\item $ 400\pi $
}

\opr{opggepvolforhold}
\abch{
	\item $\frac{32\pi}{3}$ og 288
	\item $ a^3 $
}

\opr{opgpytnamed}
\abch{
\item 15
\item 4
\item 7
\item $\sqrt{97}$
}

\opr{opggeototilen}
$c=\sqrt{5}a$

\opr{opggeopytright}
\abch{
\item Ja.
\item Nei.
}

\opr{opggeovolforholdgeneral}
$a^3$

\opr{geoopgformlvink} $ 60^\circ $

\opr{geoopgsamsv}
$ AC $ og $ DE $, $ BC $ og $ EF $, $ AB $ og $ DF $.

\opr{geofinnlen} $ EF=3 $ og $ DF=2 $.

\opr{geofinnlen2} $ AC=15 $ og $ DF=6 $.

\opr{opggeosimlar} To vinkler som er parvis like store. \selos.

\opr{opggeotop}
\abch{
\item To vinkler som er parvis like store. \selos.
\item $CD$ og $AC$, $BC$ og $EC$, $AB$ og $ED$.
\item 20.
}

\opr{opggeofinnforml} $\triangle ABC \sim \triangle ACD \sim \triangle CBD$.

\opr{opggeopytrewrite}
\abch{
\item $AC^2=BC^2+AB^2$
\item $BC^2=AC^2-AB^2$
} \os
\abchs{3}{
\item $AB^2=AC^2-BC^2$
}

\opr{opggeopyt}
\abch{
\item 10
\item 17
\item 12
\item 15
}

\opr{opggeopytright}
\abch{
\item Rettvinklet.
\item Ikke rettvinklet.
}

\grubr{opggeotrapz}
$x=3$

\grubr{opggeopytrapz}
20

\grubr{opggeopytgrub}
15 eller $3\sqrt{7}$

\grubr{opggeoeqlheight}
\abch{
	\item \selos
	\item \selos
}

\grubr{opggeo263}
$3\sqrt{10}$

\grubr{opgegeotrapes}
30

\grubr{opggeotale}
\selos

\grubr{opggeooytright2}
\abch{
\item \includegraphics[scale=0.25]{\figp{opggeooytright2a}}
\item \includegraphics[scale=0.25]{\figp{opggeooytright2b}}
\item \includegraphics[scale=0.25]{\figp{opggeooytright2c}}
}

\grubr{opggeoarforhold}
\selos

\grubr{opggeoarforhold2}
\selos

\grubr{opgegeofunktriangle}
\abch{
\item $f(x)=3x+2$, $g(x)=\frac{1}{3}+2$.
\item \includegraphics[scale=0.25]{\figp{opggeofunktriangle}} (navnene på punktene er vilkårlig valgt).
\item $\triangle ABC$ er likebeint ($AB=AC$).
\item 4
\item \selos
}

\grubr{opggeosekt}
\selos

\grubr{opggeopytar}
8, 15 og 17

\grubr{opggeopolvinksum}
$180^\circ\cdot(n-2)$

\grubr{opggeovispyt}
\selos

\grubr{opggeotrap}
\selos

\grubr{GV21D1opg12}
$30^\circ$

\grubr{opgegeotan}
\selos

\grubr{opgdoublear}
\selos

\grubr{opglikbmidtn} Se \tmen.

\grubr{opgdoublear} \selos

\grubr{opggeoperlin} \selos

\grubr{opgmedian} \selos

\grubr{opgtresirkar} \selos

\grubr{GV21D1opg12} \selos

\end{document}